\documentclass[11pt, oneside]{article}   	% use "amsart" instead of "article" for AMSLaTeX format
\usepackage{geometry}                		% See geometry.pdf to learn the layout options. There are lots.
\geometry{letterpaper}                   		% ... or a4paper or a5paper or ... 
%\geometry{landscape}                		% Activate for rotated page geometry
%\usepackage[parfill]{parskip}    		% Activate to begin paragraphs with an empty line rather than an indent
\usepackage{graphicx}				% Use pdf, png, jpg, or eps§ with pdflatex; use eps in DVI mode
								% TeX will automatically convert eps --> pdf in pdflatex		
\usepackage{amssymb}
\usepackage{amsmath}

%SetFonts

%SetFonts


\title{}
\author{}
\date{2026/03/03}							% Activate to display a given date or no date

\begin{document}
\maketitle
If $f(k)$ represents the cost of buying $1$ when $k$ are already bought, an geometric (exponential) cost scaling can be represented as:
\[
f(k) = ab^k
\]
In the code, $f(k)$ is equivalent to new ExponentialCostScaling().getCurrentCost(k).
In the code, $a=$ this.baseCost, $b=$ this.baseIncrease, $m=$ currentAmount, $n=$ buyAmount.

If m is already bought, the cost of buying n extra (C) can be represented as:
\[
C(m,n)=\sum_{k=m}^{m+n-1}f(k)
=\dfrac{ab^m(b^n-1)}{b-1}
\]
\[C(m,n)(b-1)=ab^m(b^n-1)  (b\ne1)\]
\[C(m,n)\cdot\frac{b-1}{ab^m}=b^n-1\]
\[C(m,n)\cdot\frac{b-1}{ab^m}+1=b^n\]
\[\log_b
\left( C(m,n)\cdot\frac{b-1}{ab^m}+1 \right)
=n\]
If M stands for money,
\[n=\left\lfloor
\log_b\left( M \cdot \frac{b-1}{ab^m}+1 \right)
\right\rfloor\]

%\section{}
%\subsection{}



\end{document}  